\begin{workedexample}[Worked Example: Y-factor noise figure]
A noise source with $\text{ENR} = 15\ \text{dB}$ ($=31.6$ linear) is
switched on and off at a device's input, giving a hot/cold power ratio
$Y = 8\ \text{dB}$ ($=6.31$). The noise factor is
\[ F = \frac{\text{ENR}}{Y - 1} = \frac{31.6}{6.31 - 1} = 5.95, \]
i.e.\ $\text{NF} = 7.75\ \text{dB}$. A larger $Y$ (from a lower-noise device)
gives a lower, more accurate NF reading.
\end{workedexample}

\begin{keytakeaways}
\begin{itemize}[leftmargin=1.4em,itemsep=2pt]
\item Narrower RBW lowers the analyzer noise floor and sharpens resolution but slows the sweep.
\item Y-factor: $F = \text{ENR}/(Y-1)$ from the hot/cold power ratio.
\item VNA needs SOLT/TRL calibration; TDR locates discontinuities by echo delay.
\end{itemize}
\end{keytakeaways}

\begin{exercises}
\item With $\text{ENR} = 15\ \text{dB}$ and $Y = 10\ \text{dB}$, find the noise figure.
\item What happens to a spectrum analyzer's noise floor when you narrow the RBW?
\item What does a TDR measurement locate?
\end{exercises}

\furtherreading{Pozar~\cite{pozar} (ch.~4); Steer~\cite{steer}.}
